\section{2021 Field Map：已经到哪里，还缺什么}

理解架构还需要理解它为何获得跨领域注意。到 2021 年，Transformer 已从语言生成与理解扩展到 vision、generative modeling、reinforcement learning 与结构生物学相关系统。这里不能只列应用名词；要问每个领域如何定义 token、position、mask 与 objective。

\subsection{跨领域迁移不是简单复制}

语言通常把 subword 当 token，图像可把 patch 当 token，RL 可把 state/action/reward 序列化，蛋白质可以氨基酸或 residue representation 建模。相同 self-attention 公式在不同领域处理的关系完全不同，因此 preprocessing、inductive bias 与 evaluation 决定迁移是否成立。

\lecturefigure{slide-07-where-we-are.jpg}{2021 年课堂已把 Transformer 放入 NLP、vision、RL 与生成建模的共同版图。}{官方视频 04:30。}

\begin{knowledgebox}{读图：先识别 token，再讨论架构}
在 NLP 中 token 有离散词序和语法；vision patch 还需要二维位置与局部结构；RL trajectory 包含动作因果和回报；generative image model 可能在 pixel、latent 或 codebook token 上工作。若只说“都用 Transformer”，会隐藏最关键的建模选择。读后续课程时先写四列：token、position、mask、loss。
\end{knowledgebox}

\begin{warningbox}{AlphaFold2 不是“只靠 Transformer”}
Transformer-like attention 是 AlphaFold2 的重要组成，但系统还结合多序列比对、pair representation、几何模块与结构推理。跨领域成功通常来自 attention 与领域结构的共同设计，而不是把文本模型原封不动应用到新数据。
\end{warningbox}

\subsection{Future slide 中的持久问题}

Slide 一方面想象 video understanding、finance 与 generative modeling 等应用，另一方面列出 external memory、降低计算复杂度、以及与人类目标对齐等缺失组件。这些问题之所以持久，是因为 Transformer 只定义信息混合，不自动提供长期状态、低成本计算或正确目标。

\lecturefigure{slide-08-future.jpg}{应用扩张与缺失组件必须同时阅读。}{官方视频 05:55。}

\begin{knowledgebox}{读图：右列比左列更像研究议程}
左侧应用回答“可能在哪些任务上有价值”；右侧 missing ingredients 回答“为什么架构本身还不够”。External memory 处理超出 context 的持久信息，efficiency 处理 $O(n^2)$ attention 与训练成本，alignment 处理目标、反馈和社会约束。任何应用 proposal 都应至少对应一项能力假设和一项失败模式。
\end{knowledgebox}

\teachervoice{讲者并没有把 Transformer 描述为终点，而是把应用与缺失组件并列。课堂提示：研究品味来自同时看到“能做什么”和“为什么还做不好”，而不是在成功 demo 后停止追问。}

\subsection{本章小结}

跨领域统一来自 interaction primitive，不来自任务同质化。课程的未来地图把 tokenization、memory、efficiency 和 alignment 留作后续问题，为 guest lectures 提供坐标。

